\chapter{Anforderungsanalyse}

\section{Ausgangslage}
Die aktuelle Ausgangslage des IT Supports innerhalb der Techdivision ist, dass 
der IT-Support hauptsächlich über den internen Google Chat betrieben wird.
Dies ist eine offene Chatgruppe für Mitarbeiter der Techdivsion und innerhalb 
diese können Fragen gestellt werden und sie werden von einem der Mitarbeiter 
des IT-Support Teams beantwortet. Falls es fragen sind, welche nicht über den 
Chat gelöst werden können, dann wird ein Google Meet Meeting aufgesetzt oder, falls der 
jeweilige Mitarbeiter sich im Office befindet, dann kann das auch vor Ort stattfinden.

\section{Anforderungen}
Nun wird gefordert, dass dem IT-Support Gruppenchat ein Chatbot hinzugefügt wird, 
der den Mitarbeitern des Support Teams helfen soll, indem er einfache Fragen, 
die bspw. mithilfe von Dokumentationen innerhalb des internen Confluences 
beatnwortbar wären, beantworten soll. \par
Im Folgenden werden die Rahmenbedingungen, die der Chat erfüllen sollte 
vorgestellt.

\subsection{A1 Der Bot soll selbst gehostet werden}
Das größte Problem, welches beim Betreiben des Chatbots beachtet, ist die 
DSGVO. Datenschutz ist ein großes Thema bei Unternehmen und vor allem, bei 
Unternehmen die mit Kundendaten arbeiten müssen oder auch Betriebsgehemnisse,
die nicht an die Öffentlichkeit weiter gegeben werden sollen. 
Die Techdivision ist ein IT-Dienstleister und somit ist dies auch ein
Bereich der bei der Konzeption und Entwicklung dieses Chatbots beachtet werden 
muss. Aufgrund dieser Sorgen wird gefordert, dass der Chatbot von der 
Techdivision gehostet wird, was bedeutet es dürfen keine APIs von externen LLM-Systemen 
hierfür verwendet
werden, die von anderen Unternehmen gehostet werden, wie beispielsweise
die ChatGPT API.
\subsection{A2 Der Bot soll in einer Chatgruppe laufen}
Auch, wenn die LLMs heutzutage weit entwickelt sind, können diese die
falschen Informationen weiter geben, oder komplett neue halluzinieren
und diese dann dem Supportsuchenden weiter geben. Es könnte auch sein, dass
der Bot dem Mitarbeiter überhaupt nicht helfen kann. Damit in diesen Fällen
ein Supportmitarbeiter eingreifen kann, sollte der Bot in einer Chatgruppe 
laufen, von der ein Mitarbeiter des Support Teams ebenfalls Teil ist.
\subsection{A3 Der Bot soll Unternehmensdaten verwenden}
Da die Bots nicht von Haus aus Zugriff auf spezifische Informationen 
der Techdivion Zugriff haben, soll der Bot auf Unternehmensdaten 
zugreifen können. Dadurch kann der Bot spezifischer Hilfe leisten und 
nicht nur allgemeine Hinweise geben, die im Internet zu finden sind.
\subsection{A4 Der Bot soll mit der Zeit auf neues Wissen zugreifen können}
Da sich die Unternehmensdaten mit der Zeit ändern und auch neue 
Datenquellen hinzukommen können, sollte der Bot in der Lage sein 
mit der Zeit auf neues Wissen zuzugreifen.
\subsection{A5 Der Bot soll auf Deutsch arbeiten}
Der Großteil der Techdivion ist deutschsprachig. Die meisten Teile von 
Datenquellen, wie beispielsweise Confluence sind ebenfalls auf Deutsch. 
Es gibt ein paar englischsprachige Mitarbeiter, daher wäre ein Bilingualer 
Bot der beste Fall, da aber die meisten Deutsch sind, wäre Deutsch die bessere
Option, falls es nur eine Sprache sein kann.
\subsection{A6 Der Bot soll Anfragedaten speichern können}
Dies ist eine optionale Anfrage des IT-Supports.
Der Bott soll Anfragedaten speichern können, damit eingesehen werden kann,
ob zu manchen Themen viele Fragen gestellt werden oder manche Personen 
besonders viele Fragen stellen. So könnte, falls notwendig zu einem bestimmten
Thema ein Workshop angeboten werden oder es könnte ein Kurs mit einer 
Person geben, in der diese Person alles, was sie will erklärt kriegt.
